% !TeX root = ../../main.tex
\subsection{原級のつくり方}\vspace{-0.2em}
    原級は\fitblankbf[]{変化しない}。\fitblankbf[]{as}と\fitblankbf[]{as}で原形をはさむ。

    \begin{definitionbox}{\textbf{原級のつくり方}}
        \begin{enumerate}
            \item 原形(\fitblankbf[]{原級})はそのまま使う
            \compareexrow{\comparecelllong{tall}{as tall as}}{\comparecelllong{fast}{as fast as}}
            \compareexrow{\comparecelllong{beautiful}{as beautiful as}}{\comparecelllong{good}{as good as}}
            \vspace{0.65em}
            \item 肯定形は \fitblankbf[]{as} + 原級 + \fitblankbf[]{as} ～
            \vspace{0.65em}
            \item 否定形は \fitblankbf[]{not as} + 原級 + \fitblankbf[]{as} ～
        \end{enumerate}
    \end{definitionbox}
    \vspace{-0.8em}

\subsection{原級を使った構文}\vspace{-0.4em}
    \begin{theorembox}{\textbf{as＋原級＋as}}
        \theoremline{as ＋ 原級 ＋ as …}{「～と\fitblankbf[]{同じくらい}…」}
    \end{theorembox}
    \vspace{0.5em}

    \begin{theorembox}{\textbf{not as＋原級＋as}}
        \theoremline{not as ＋ 原級 ＋ as …}{「～\fitblankbf[]{ほど}…\fitblankbf[]{でない}」}
    \end{theorembox}
    \vspace{0.5em}

    \begin{theorembox}{\textbf{times as＋原級＋as}}
        \theoremline{times as ＋ 原級 ＋ as …}{「～の□\fitblankbf[]{倍}…」}
    \end{theorembox}
    \vspace{0.35em}

    ※『2倍』のときは\fitblankbf[]{twice}(＝ two times)。
    『半分、$\dfrac{1}{2}$』のときは\fitblankbf[]{half}が使われることもある。
    \vspace{0.15em}

    \begin{theorembox}{\textbf{as＋原級＋as＋人＋can [ possible ]}}
        as ＋ 原級 ＋ as ＋ 人 ＋ (\fitblankbf[]{can})［ possible ］ …
        \theoremmeaning{「\fitblankbf[]{できるだけ}～」}
    \end{theorembox}
    \vspace{0.15em}

\newpage

\subsection{練習問題}
    \begin{enumerate}[leftmargin=*, nosep, itemsep=2.0em]
        \item 彼は私と同じくらい速く泳げる。
        \dialogueblank{He can swim as fast as me.}
        \item My brother is as tall as my father.
        \dialogueblank{私の兄は父と同じくらい背が高い。}
        \item この部屋はあの部屋ほど広くない。
        \dialogueblank{This room is not as large as that one.}
        \item She doesn't run as fast as her sister.
        \dialogueblank{彼女は姉ほど速く走らない。}
        \item 東京は札幌の2倍大きい。
        \dialogueblank{Tokyo is twice as large as Sapporo.}
        \item This problem is three times as difficult as that one.
        \dialogueblank{この問題はあの問題の3倍難しい。}
        \item できるだけ大きな声で話してください。
        \dialogueblank{Please speak as loudly as you can.}
    \end{enumerate}

\newpage
